% ============================================================
% 山东大学论文模板配置
% ============================================================

% ------------------- LaTeX3 key-value 配置 -------------------
\RequirePackage{expl3}
\RequirePackage{l3keys2e}

% 使用 LaTeX3 语法模式
\ExplSyntaxOn

% 声明存储配置值的 token list 变量
\tl_new:N \l__sdu_title_tl
\tl_new:N \l__sdu_author_tl
\tl_new:N \l__sdu_studentid_tl
\tl_new:N \l__sdu_school_tl
\tl_new:N \l__sdu_major_tl
\tl_new:N \l__sdu_supervisor_tl
\tl_new:N \l__sdu_year_tl
\tl_new:N \l__sdu_month_tl

% 定义配置键
\keys_define:nn { sdu } {
  title       .tl_set:N = \l__sdu_title_tl,
  author      .tl_set:N = \l__sdu_author_tl,
  studentId   .tl_set:N = \l__sdu_studentid_tl,
  school      .tl_set:N = \l__sdu_school_tl,
  major       .tl_set:N = \l__sdu_major_tl,
  supervisor  .tl_set:N = \l__sdu_supervisor_tl,
  year        .tl_set:N = \l__sdu_year_tl,
  month       .tl_set:N = \l__sdu_month_tl,
}

% SDUSetup 命令定义
\NewDocumentCommand \SDUSetup { m } {
  \keys_set:nn { sdu } { #1 }
}

% Getter 命令：供模板内部和封面使用
\NewDocumentCommand \GetTitle { } { \l__sdu_title_tl }
\NewDocumentCommand \GetAuthor { } { \l__sdu_author_tl }
\NewDocumentCommand \GetStudentId { } { \l__sdu_studentid_tl }
\NewDocumentCommand \GetSchool { } { \l__sdu_school_tl }
\NewDocumentCommand \GetMajor { } { \l__sdu_major_tl }
\NewDocumentCommand \GetSupervisor { } { \l__sdu_supervisor_tl }
\NewDocumentCommand \GetYear { } { \l__sdu_year_tl }
\NewDocumentCommand \GetMonth { } { \l__sdu_month_tl }

\ExplSyntaxOff

% ------------------- LaTeX3 key-value 配置 -------------------

% 导入宏包
\usepackage{amsmath}
\usepackage{amsthm}
\usepackage{amsfonts}
\usepackage{amssymb}
\usepackage{unicode-math}
\usepackage{xeCJK}
\usepackage{xcolor}
\usepackage{geometry}
\usepackage{float}
\usepackage{fancyhdr}
\usepackage{setspace}
\usepackage{bookmark}
\usepackage{booktabs}
\usepackage{graphicx}
\usepackage{caption}
\usepackage{listings}
% algorithm2e 必须在其他宏包之后加载，避免冲突
\usepackage{tocloft}
\usepackage{subfig}
\usepackage[
    backend=biber,
    style=gb7714-2015,
    gbnamefmt=givenahead
]{biblatex}
\usepackage{hyperref}
\usepackage{etoolbox}
\usepackage{tabularx}

\usepackage{algorithmicx}
\usepackage{algorithm}
\usepackage{algpseudocode}
% -------------------算法设置-------------------
% \Return 已由 algorithmicx 预定义，直接使用即可
\AtBeginEnvironment{algorithm}{\linespread{0.85}\selectfont}


% -------------------文本选项-------------------
% 让已定义的中文字体失效
\let\songti\relax
\let\heiti\relax
\let\kaiti\relax

% 字体文件检测命令（已简化，不检查本地字体文件）
\def\IfFontExists#1#2#3{#3}

% 中文字体配置
% 使用 ctex 自动检测的中文字体，或使用开源 Fandol 字体
\newCJKfontfamily\songti{FandolSong}
\newCJKfontfamily\heiti{FandolHei}
\newCJKfontfamily\kaiti{FandolKai}

% 后续局部定义字体，后面必须带{}
\newcommand{\song}[1]{{\songti{#1}}}
\newcommand{\hei}[1]{{\heiti{#1}}}
\newcommand{\kai}[1]{{\kaiti{#1}}}

% 后续局部变形处理，对于每种字体定义了加粗、倾斜和加粗倾斜，后面必须带{}
\newcommand{\bfsong}[1]{{\songti\textbf{#1}}}
\newcommand{\bfhei}[1]{{\heiti\textbf{#1}}}
\newcommand{\bfkai}[1]{{\kaiti\textbf{#1}}}
\newcommand{\itsong}[1]{{\songti\textit{#1}}}
\newcommand{\ithei}[1]{{\heiti\textit{#1}}}
\newcommand{\itkai}[1]{{\kaiti\textit{#1}}}
\newcommand{\bfitsong}[1]{{\songti\textbf{\textit{#1}}}}
\newcommand{\bfithei}[1]{{\heiti\textbf{\textit{#1}}}}
\newcommand{\bfitkai}[1]{{\kaiti\textbf{\textit{#1}}}}

% 后续全局变形处理，对于每种字体定义了加粗、倾斜和加粗倾斜，后面不能带{}
\newcommand{\allbfsong}{\songti\bfseries}
\newcommand{\allbfhei}{\heiti\bfseries}
\newcommand{\allbfkai}{\kaiti\bfseries}
\newcommand{\allitsong}{\songti\itshape}
\newcommand{\allithei}{\heiti\itshape}
\newcommand{\allitkai}{\kaiti\itshape}
\newcommand{\allbfitsong}{\songti\bfseries\itshape}
\newcommand{\allbfithei}{\heiti\bfseries\itshape}
\newcommand{\allbfitkai}{\kaiti\bfseries\itshape}

% 设置英文字体（使用开源替代字体）
\setmainfont{TeX Gyre Termes}
\setsansfont{TeX Gyre Heros}
\setmonofont{TeX Gyre Cursor}

% 设置数学模式下的字体
\setmathfont{Latin Modern Math}

% 全局默认使用小四号宋体
\AtBeginDocument{\zihao{-4}\songti}

% 全局默认使用1.5倍行距
\setstretch{1.5}
% -------------------文本选项-------------------


% -------------------页面选项-------------------
% 页边距设置
\geometry{
    a4paper,
    left=3cm,
    right=3cm,
    top=2.5cm,
    bottom=2.5cm,
}

% 页面顶部对齐，允许底部有空白
\raggedbottom
% -------------------页面选项-------------------


% -------------------章节标题选项-------------------
% 最多显示三级标题
\setcounter{secnumdepth}{3}

% 章节标题格式与前后距离
\ctexset{
    chapter={
      format=\centering\zihao{3}\allbfhei,
      name={第,章},
      beforeskip=21.6pt,
      afterskip=18pt,
      fixskip=true,
     },
    section={
      format= \zihao{4}\allbfhei,
      name={},
      beforeskip=18pt,
      afterskip=18pt,
      fixskip=true,
     },
    subsection={
            format=\zihao{-4}\allbfhei,
            name={},
            beforeskip=18pt,
            afterskip=18pt,
            fixskip=true,
        }
}
% -------------------章节标题选项-------------------


% -------------------交叉引用选项-------------------
% 自定义颜色
\definecolor{linkdarkblue}{rgb}{0,0.08,0.45}

\hypersetup{
    colorlinks=true,       % 启用彩色链接
    linkcolor=black,       % 文本内部链接颜色
    citecolor=black,       % 引用文献的链接颜色
    filecolor=black,       % 文件链接颜色
    urlcolor=linkdarkblue, % URL链接颜色
}

% PDF属性选项
\hypersetup{
    bookmarksnumbered=true,
    pdfstartview=FitH,
}

% 自定义引用命令
\newcommand{\equref}[1]{式\ \eqref{#1}\ }
\newcommand{\tabref}[1]{表\ \ref{#1}\ }
\newcommand{\figref}[1]{图\ \ref{#1}\ }
\newcommand{\subfigref}[1]{{图\ \subref*{#1}\ }}
% biblatex 引用命令（兼容不同风格）
\let\citing\supercite  % 上标数字引用 [1]
\let\citex\citep       % 括号引用 (作者, 年份)
% -------------------交叉引用选项-------------------


% -------------------图表标题选项-------------------
% 图片存放路径
\graphicspath{{figures/}}

% 自定义标题样式
\DeclareCaptionFont{mybfsong}{\allbfsong}
\DeclareCaptionFont{mywuhao}{\zihao{5}}
\DeclareSubrefFormat{mysubref}{#1(#2)}

\captionsetup[figure]{
    position=below,
    font={mywuhao,stretch=1.0},
    labelfont=mybfsong,
    textfont=mybfsong,
    labelsep=quad,
    margin=40pt,
    belowskip=-4pt,
}
\captionsetup[table]{
    position=above,
    font={mywuhao,stretch=1.0},
    labelfont=mybfsong,
    textfont=mybfsong,
    labelsep=quad,
    margin=40pt,
    belowskip=4pt,
}
\captionsetup[subfloat]{
    position=below,
    font={mywuhao,stretch=1.0},
    labelfont=mybfsong,
    textfont=mybfsong,
    subrefformat=mysubref,
}
% -------------------图表标题选项-------------------


% -------------------数学选项-----------------------
\newcommand{\mye}{\ensuremath{\mathrm{e}}}
\newcommand{\myi}{\ensuremath{\mathrm{i}}}
\newcommand{\myj}{\ensuremath{\mathrm{j}}}
% -------------------数学选项-----------------------


\lstset{
  belowskip=0pt,
  language=C++,
  basicstyle=\ttfamily\small,
  numbers=left,
  numberstyle=\tiny\color{gray},
  numbersep=5pt,
  breaklines=true,
  frame=single,
  rulecolor=\color{gray},
  framexleftmargin=1em,
  xleftmargin=2em,
  captionpos=b,
  showstringspaces=false,
  tabsize=4,
  columns=flexible,
  backgroundcolor=\color{gray!5},
  morekeywords={constexpr, nullptr, Ciphertext, Plaintext, vector, load_from_file, save_to_file}
}
